\section*{अध्याय \ref{ch:b2:basins-compared} --- अवसादी द्रोणियों की तुलना}

\begin{solution}{exo:b2:basins-compared:1}
(क) अंतःक्रेटॉनिक: सैकड़ों मिलियन वर्षों तक स्पंदों में धीमा अवधसान, बड़े
भ्रंशन के बिना; उसके कारण पर बहस है। (ख) दरार: खिंचाव के दौरान भूपर्पटी का यांत्रिक पतलापन,
फिर स्थलमंडल के ठंडा होने पर तापीय संकुचन। (ग) निष्क्रिय किनारा: दरार-चरण से विरासत में मिला
तापीय अवधसान, अवसाद की मोटी कील के भार से प्रवर्धित। (घ) और (ङ)
अग्रभूमि: अधिक्षेप-पट्टी, आल्प्स और हिमालय, के भार से नीचे झुका स्थलमंडल।
\end{solution}

\begin{solution}{exo:b2:basins-compared:2}
(क) $Y = 3000\times(3330 - 2300)/(3330 - 1030) = 3000\times1030/2300 = \qty{1343}{m}$।
(ख) $1343 + 200 = \qty{1543}{m}$। (ग) अपने भार के बिना, आधारशैल \qty{1343}{m} पर होता:
मोटाई के $1 - 1343/3000 = \qty{55}{\percent}$ के लिए जगह भार बनाता है।
\end{solution}

\begin{solution}{exo:b2:basins-compared:3}
(क) $\beta = 35/25 = 1.4$। (ख) $125/1.4 = \qty{89.3}{km}$। (ग) $S_i = 4334\times(1 - 1/1.4) =
\qty{1238}{m}$।
\end{solution}

\begin{solution}{exo:b2:basins-compared:4}
(क) $\tau = (1.25\times10^{5})^2/(\pi^2\times8\times10^{-7}) = \qty{1.979e15}{s}$, यानी
$1.979\times10^{15}/3.156\times10^{13} = \qty{62.7}{Myr}$। (ख) $62.7\times0.8 = \qty{50.2}{Myr}$।
(ग) $62.7\times(100/125)^2 = \qty{40.1}{Myr}$। (घ) $1 - e^{-1} = \qty{63}{\percent}$; $1 -
e^{-3} = \qty{95}{\percent}$।
\end{solution}

\begin{solution}{exo:b2:basins-compared:5}
(क) $0.63\,e^{-1000/1960} = 0.378$; $0.63\,e^{-2000/1960} = 0.227$। (ख) $T_s = 1000 -
0.63\times1960\times(0.6004 - 0.3605) = \qty{704}{m}$। (ग) $m = 2700\times704 + 1030\times296 =
\qty{2.206e6}{kg.m^{-2}}$; $\bar\rho = \qty{2206}{kg.m^{-3}}$।
\end{solution}

\begin{solution}{exo:b2:basins-compared:6}
(क) $z_2' \leftarrow 704 + 1234.8\,(1 - e^{-z_2'/1960})$: 1077, 1226, 1278, 1296, 1301, 1303,
1304; परत \qty{1304}{m} मोटी है। (ख) $(2700\times704 + 1030\times600)/1304 =
\qty{1932}{kg.m^{-3}}$। (ग) \qty{304}{m}, यानी \qty{30}{\percent}, अधिक: वह जल जो संपीड़न ने
निचोड़ दिया है।
\end{solution}

\begin{solution}{exo:b2:basins-compared:7}
(क) $1 - e^{-100/62.7} = 0.7971$। $\beta = 1.5$ के लिए: $3207\times0.4135\times0.7971 =
\qty{1057}{m}$; $\beta = 2$ के लिए: $3207\times0.6366\times0.7971 = \qty{1627}{m}$। (ख) हाँ: 1100 से
\qty{1400}{m} इनके बीच है। (ग) $\frac{\beta}{\pi}\sin\frac{\pi}{\beta} = 1250/(3207\times0.7971)
= 0.4890$, अतः $\beta = 1.6 + 0.1\times(0.4890 - 0.4705)/(0.5204 - 0.4705) = 1.64$, जो द्रोणी के लिए अनुमानित
50 से \qty{100}{\percent} के खिंचाव के भीतर है।
\end{solution}

\begin{solution}{exo:b2:basins-compared:8}
(क) हर में $\rho_w$ का स्थान भराव का घनत्व लेता है। (ख) अंश
$125\,000\times(1/3)\times74.69 = \num{3.112e6}$ है; $S_i = 3.112\times10^{6}/(3184.4 - 2300) =
\qty{3519}{m}$। (ग) जल के नीचे, $3.112\times10^{6}/2154.4 = \qty{1445}{m}$: अवसाद से भरी
द्रोणी 2.44 गुना गहरी है।
\end{solution}

\begin{solution}{exo:b2:basins-compared:9}
(क) $2171 - 100\times3330/2300 = 2171 - 145 = \qty{2026}{m}$। (ख) $Y$ मीटर-दर-मीटर
$W_d$ के साथ बदलता है: 1971 से \qty{2371}{m} तक। (ग) जल की गहराई, $\pm\qty{200}{m}$, बनाम
समुद्र-स्तर के लिए \qty{145}{m}।
\end{solution}

\begin{solution}{exo:b2:basins-compared:10}
(क) $S_i$ का चिह्न उसके अंश के कोष्ठक का चिह्न है, क्योंकि हर, $3184.4 - 1030$,
धनात्मक है; वह ऋणात्मक, यानी उत्थान, तब होता है जब कोष्ठक ऋणात्मक हो। (ख) $(\rho_m - \rho_c)t_c/a =
\alpha T_1\rho_m/2$ से $t_c = a\,\alpha T_1\rho_m/(2(\rho_m - \rho_c)) = 125\times0.04372\times
3330/1060 = \qty{17.17}{km}$ मिलता है। (ग) $t_c = 17.17/(1 - 0.04372\times17.17/250) = \qty{17.22}{km}$।
(घ) जहाँ खिंचाव से पहले भूपर्पटी लगभग \qty{17}{km} से पतली हो: महासागरीय स्थलमंडल में,
या किसी पहले के दरार-निर्माण से पहले ही पतली हुई महाद्वीपीय भूपर्पटी में।
\end{solution}

\begin{solution}{exo:b2:basins-compared:11}
(क) खिंचाव से पहले, स्तंभ में द्रव्यमान $M = \rho_c t_c(1 - \alpha T_1
t_c/(2a)) + \rho_m(a - t_c)(1 - \alpha T_1(a + t_c)/(2a))$ का स्थलमंडल है; बाद में, $\rho_w S_i + M/\beta +
\rho_m(1 - \alpha T_1)(a - S_i - a/\beta)$। उन्हें बराबर करने पर, $S_i[\rho_m(1 - \alpha T_1) - \rho_w] =
(1 - 1/\beta)[\rho_m(1 - \alpha T_1)a - M]$, और $M$ का विस्तार करने पर, $\rho_m(1 - \alpha T_1)a - M =
(\rho_m - \rho_c)t_c(1 - \alpha T_1 t_c/(2a)) - \rho_m\alpha T_1 a/2$: यही सूत्र है। (ख) अंतिम प्रोफ़ाइल $T_1 z/a$ है; प्रारंभिक प्रोफ़ाइल उससे
$T_1 z(\beta - 1)/a$ अधिक है $a/\beta$ के ऊपर, और $T_1(1 - z/a)$ अधिक है नीचे। समाकल $T_1
a(\beta - 1)/(2\beta^2)$ और $T_1 a(\beta - 1)^2/(2\beta^2)$ हैं, जिनका योग $T_1 a(1 - 1/\beta)/2$ है।
इसका $\alpha$ गुना संकुचन, जल से भारित, व्यंजक देता है। (ग) $3.28\times10^{-5}
\times3330\times1333\times125\,000\times0.5/(2\times2300) = \qty{1978}{m}$, बनाम $3207\times
2/\pi = \qty{2042}{m}$: पहला पद अंतिम अवधसान को \qty{64}{m},
\qty{3}{\percent}, अधिक बताता है।
\end{solution}

\begin{solution}{exo:b2:basins-compared:12}
(क) $3207\times0.6366\times(1 - e^{-150/62.7}) = 2042\times0.9086 = \qty{1855}{m}$। (ख) $1855\times
2300/(3330 - 2400) = \qty{4588}{m}$। (ग) 6.9 से \qty{13.8}{km}। (घ) प्रतिरूप कई ऐसी बातें
छोड़ देता है जो द्रोणी को गहरा कर सकती हैं: भूपर्पटी से अधिक खिंचा प्रावार-स्थलमंडल, अवसाद के भार की
क्षेत्रीय (आनमनी) प्रतिक्रिया, दरार-निर्माण के दौरान किनारे का अस्थायी उत्थान
जिसके बाद धँसाव हुआ, या अवसाद की इतनी बड़ी आपूर्ति कि द्रोणी आवश्यकता से अधिक भर गई और शेल्फ़ पर
समुद्र-तल से ऊपर तक भारित हुई।
\end{solution}

\begin{solution}{pb:b2:basins-compared:1}
\textbf{1.} शीर्ष पर $0.49\,e^{-4420/3704} = 0.149$; आधार पर $0.49\,e^{-6030/3704} = 0.096$।
\textbf{2.} $T_s = 1610 - 0.49\times3704\times(0.3032 - 0.1963) = \qty{1416}{m}$।
\textbf{3.} शेल: $1110 - 0.63\times1960\times(e^{-3310/1960} - e^{-4420/1960}) = 1110 - 1234.8
\times(0.1847 - 0.1049) = \qty{1011}{m}$।
\textbf{4.} शेलयुक्त रेत: $2680\times1004 + 1030\times696 = \qty{3.408e6}{kg.m^{-2}}$; चॉक:
$2710\times1409 + 1030\times201 = \qty{4.025e6}{kg.m^{-2}}$; शेल: $2700\times1011 + 1030\times99 =
\qty{2.832e6}{kg.m^{-2}}$; रेत: $2650\times1416 + 1030\times194 = \qty{3.952e6}{kg.m^{-2}}$। कुल
\qty{1.4217e7}{kg.m^{-2}}, $\bar\rho_s = 14\,216\,910/6030 = \qty{2358}{kg.m^{-3}}$।
\textbf{5.} $z_2' \leftarrow 1416 + 1815\,(1 - e^{-z_2'/3704})$: 1993, 2171, 2221, 2235, 2238, 2239,
2239: \qty{2239}{m}।
\textbf{6.} $(2650\times1416 + 1030\times823)/2239 = \qty{2055}{kg.m^{-3}}$।
\textbf{7.} $Y = 2239\times(3330 - 2055)/2300 + 200 = 1241 + 200 = \qty{1441}{m}$।
\textbf{8.} $Y = 6030\times(3330 - 2358)/2300 + 100 = 2548 + 100 = \qty{2648}{m}$। इकाइयाँ
अपनी वर्तमान गहराइयों पर मापी गई हैं: वे पहले से उतनी ही संपीडित हैं जितनी आज हैं।
\textbf{9.} $(2171 - 1441)/50 = \qty{14.6}{m}$ प्रति मिलियन वर्ष, 150 से \qty{100}{Ma} तक;
$(2453 - 2171)/40 = \qty{7.05}{m}$ 100 से \qty{60}{Ma} तक; $(2648 - 2453)/60 = \qty{3.25}{m}$
तब से: धँसाव मंद होता जाता है।
\textbf{10.} $6030 + 100 = \qty{6130}{m}$, विवर्तनिक अवधसान का 2.31 गुना।
\textbf{11.} $(6130 - 2648)/6130 = \qty{57}{\percent}$।
\textbf{12.} $2171 - 100\times3330/2300 = \qty{2026}{m}$।
\textbf{13.} दरार-चरण के अंत का बिंदु, \qty{150}{Ma}: $S_i = \qty{1441}{m}$, यदि
दरार-निर्माण के \qty{15}{Myr} के दौरान कोई तापीय अवधसान न हुआ हो।
\textbf{14.} $\beta = 1/(1 - 1441/4334) = 1/0.6675 = 1.50$।
\textbf{15.} $35/1.50 = \qty{23.3}{km}$ भूपर्पटी और $125/1.50 = \qty{83.3}{km}$ स्थलमंडल।
\textbf{16.} $(1.50/\pi)\sin(\pi/1.50) = 0.4775\times0.8660 = 0.4135$; $3207\times0.4135 =
\qty{1326}{m}$।
\textbf{17.} $S_i = 4334/3 = \qty{1445}{m}$। \qty{100}{Ma} पर, $t = \qty{50}{Myr}$: $1326\times(1 -
e^{-50/62.7}) = \qty{729}{m}$, $Y = \qty{2173}{m}$; \qty{60}{Ma} पर: $1326\times0.7620 =
\qty{1010}{m}$, $Y = \qty{2455}{m}$; आज: $1326\times0.9086 = \qty{1205}{m}$, $Y = \qty{2650}{m}$।
\textbf{18.} प्रतिरूप हर बिंदु पर \qty{2}{m} गहरा है: 2171 के विरुद्ध 2173, 2453 के विरुद्ध 2455,
2648 के विरुद्ध 2650।
\textbf{19.} पुरा-जल-गहराइयाँ केवल सौ मीटर या अधिक की परिशुद्धता तक ज्ञात; अतीत के समुद्र-स्तर
आज से भिन्न; \qty{15}{Myr} तक चला दरार-निर्माण, जिसके दौरान ऊष्मा का एक भाग
पहले ही खो गया; स्थानीय समस्थिति के बजाय आनमन; भूपर्पटी से अधिक या कम खिंचा प्रावार;
दरार-निर्माण से पहले \qty{35}{km} से भिन्न मोटाई वाली भूपर्पटी।
\textbf{20.} $1326\times(e^{-150/62.7} - e^{-200/62.7}) = \qty{67}{m}$।
\textbf{21.} $66.6\times2300/(3330 - 2400) = \qty{165}{m}$।
\textbf{22.} $1 - e^{-150/62.7} = \qty{91}{\percent}$।
\textbf{23.} दरार-चरण के अंत के $\tau\ln2 = \qty{43.5}{Myr}$ बाद: \qty{106.5}{Ma} पर।
\textbf{24.} $\beta = \textbf{1.50}$: भूपर्पटी अपनी चौड़ाई के डेढ़ गुना तक खींची गई थी।
\end{solution}
